\section*{Chapter \ref{ch:m1:auctions} --- Auctions}

\begin{solution}{exo:m1:auctions:1}
At 19.95 and 20.00: $D = 1\,200$, $S = 200$, executable 200. At 20.05: $D =
700$, $S = 800$, executable 700. At 20.10: $D = 400$, $S = 1\,100$,
executable 400. The auction uncrosses at 20.05 for 700 shares.
\end{solution}

\begin{solution}{exo:m1:auctions:2}
Surplus $700 - 800 = -100$: 100 shares remain to sell at 20.05. Both buy
orders at 20.10 and 20.05 are filled in full; the sell order at 19.95 in full
and the one at 20.05 for 500 of its 600. The orders at 20.00 (buy) and 20.10
(sell) do not trade.
\end{solution}

\begin{solution}{exo:m1:auctions:3}
On Nasdaq, freely: MOC entry closes at 15:55. On the New York Stock Exchange
only if a regulatory \emph{sell} imbalance was published at 15:50 in that
symbol: after the cut-off closing orders are accepted only on the side that
offsets the published imbalance.
\end{solution}

\begin{solution}{exo:m1:auctions:4}
Every price from 50.00 to 50.20 executes 1\,000 shares with zero surplus:
rules 1 to 3 do not decide. Rule 4 gives the reference price, 50.12. With a
reference of 50.40, outside the range, the closest candidate is 50.20. Two
parties willing to trade anywhere in a 20-cent range are priced by the last
trade, not by either of them.
\end{solution}

\begin{solution}{exo:m1:auctions:5}
$300\,000/(25\,000 + 40\,000) = 4.6$ ticks, about 5 cents, that is
$4.6/8\,000 = \qty{5.8}{\bp}$.
\end{solution}

\begin{solution}{exo:m1:auctions:6}
At the close: $0.12\% \times 50 = \$60\,000$ above the current price. Now:
$0.08\% \times 50 = \$40\,000$ of impact, \emph{plus} the risk that the close
differs from the price obtained, which the fund cannot justify: its
benchmark is the close, and a purchase at the close has zero \omterm{def:m1:the-buy-side:benchmark}{tracking error}
whatever the premium. The \$20\,000 is the price of that certainty, paid by
the fund's investors and invisible in its tracking statistics.
\end{solution}

\begin{solution}{exo:m1:auctions:7}
With 400 orders an imbalance of 50\% (about 45\,000 shares) moves the close
by 1.83 ticks: 4.1 ticks per 100\,000 shares. With 1\,600 orders the same
50\% is about 182\,000 shares and moves it 1.87 ticks: 1.0 tick per 100\,000
shares. Four times as many orders means four times the density of supply and
demand per tick, hence a quarter of the move per share, as
\cref{prop:m1:auctions:depth} predicts; an imbalance that is a fixed
\emph{fraction} of the auction moves the price by the same amount in both.
\end{solution}

\begin{solution}{exo:m1:auctions:8}
(i) At the \omterm{def:m1:auctions:call}{uncrossing price} orders are filled in time priority and the
surplus side is only partly filled: an order entered at 15:55 at exactly the
\omterm{def:m1:auctions:moc}{closing price} is behind everyone earlier and often receives nothing. (ii) The
\omterm{def:m1:auctions:imbalance}{indicative price} at 15:55 is itself moved by the strategy's order and by the
offsetting interest it attracts; when the close moves \emph{through} the
limit, it is typically because better-informed interest arrived afterwards,
so the fills obtained are adversely selected. Correcting this requires the
auction's paired and imbalance quantities at the uncrossing, the order's
position in the time queue at its limit, and an estimate of the strategy's
own impact from \cref{prop:m1:auctions:depth}.
\end{solution}

\begin{solution}{pb:m1:auctions:1}
\textbf{1.} $D$: 400\,000 at 39.90; 250\,000 at 39.95 and 40.00; 130\,000 at
40.05 and 40.10; 50\,000 from 40.15. $S$: 60\,000; 160\,000; 160\,000;
300\,000; 300\,000; 500\,000 at 40.15; 750\,000 at 40.30; 1\,050\,000 at
40.50.
\textbf{2.} Executable: 60, 160, 160, 130, 130, 50, 50, 50 (thousands). The
maximum, 160\,000, occurs at 39.95 and 40.00 with the same surplus; the
surplus is to buy at both, so rule 3 takes the higher: 40.00.
\textbf{3.} $250\,000 - 160\,000 = 90\,000$ to buy.
\textbf{4.} The buy order at 40.00 receives 30\,000 of its 120\,000.
\textbf{5.} Every value of $D$ rises by 500\,000: 900, 750, 750, 630, 630,
550, 550, 550 (thousands).
\textbf{6.} Executable: 60, 160, 160, 300, 300, 500, 550, 550. The maximum
550\,000 occurs at 40.30 and 40.50; the surplus is 200\,000 to sell at 40.30
and 500\,000 at 40.50: rule 2 gives 40.30.
\textbf{7.} 30 cents, 75 basis points.
\textbf{8.} At 40.00: $750\,000 - 160\,000 = 590\,000$ shares to buy.
\textbf{9.} $0.30 \times 500\,000 = \$150\,000$.
\textbf{10.} $S$ becomes 450\,000 at 40.10, 650\,000 at 40.15, 850\,000 at
40.20, 1\,100\,000 at 40.30. Executable 550\,000 from 40.15 upward; the
smallest surplus, 100\,000 to sell, is at 40.15: the close is 40.15 for
550\,000 shares.
\textbf{11.} $0.15 \times 500\,000 = \$75\,000$.
\textbf{12.} It sold at 40.15 and expects to buy back at 40.05: $0.10 \times
150\,000 = \$15\,000$.
\textbf{13.} It is short 150\,000 shares overnight in a stock that has just
entered an index: \$6 million of single-name risk, including any news. It can
hedge the market component with index futures, or have bought the shares
beforehand, during the days between announcement and inclusion --- in which
case its sale at the close completes an inventory trade.
\textbf{14.} Between 40.00 and 40.15 supply rises by 490\,000 and demand
falls by 200\,000: about 46\,000 shares per tick. The proposition predicts
$500\,000/46\,000 = 10.9$ ticks against the 15 observed: the book is lumpy,
and the index order exhausted the dense region near 40.00.
\textbf{15.} A lower average price, by anticipating their own demand; at the
cost of \omterm{def:m1:the-buy-side:benchmark}{tracking error} against an index that includes the stock only from
tonight's close, and of revealing their purchases to those who trade ahead
of them anyway.
\textbf{16.} The index funds' investors, through a \omterm{def:m1:auctions:moc}{closing price} 15 cents
above the pre-auction market: the providers' \$15\,000 is part of it, the
sellers of the limit orders hit on the way up receive the rest.
\textbf{17.} To push the \omterm{def:m1:auctions:imbalance}{indicative price} up and induce others to enter sell
interest or to give up bidding, then withdraw: it never intended to trade.
This is spoofing, a form of market manipulation; at the close it also marks
the \omterm{def:m1:auctions:moc}{closing price}.
\textbf{18.} So that the published imbalance is a credible quantity that can
only shrink: late orders that could add to it, or cancellations that could
resurrect it, would make the information worthless and the \omterm{def:m1:auctions:moc}{closing price}
easy to manipulate.
\textbf{19.} \textbf{40.15 for 550\,000 shares.}
\textbf{20.} The certainty of owning the stock at exactly the price its
benchmark uses.
\end{solution}

\begin{solution}{iq:m1:auctions:1}
Orders accumulate without trading. At the close the exchange finds the price
that maximises the matched volume; ties are broken by the smallest imbalance,
then by the side of the imbalance, then by proximity to a reference price.
All matched orders trade at that one price, with market orders and
better-priced limits first and time priority at the clearing price.
\iqlookfor{maximum volume first, and a sensible tie-break sequence.}
\end{solution}

\begin{solution}{iq:m1:auctions:2}
Index funds and benchmarked managers are valued at the close and avoid
\omterm{def:m1:the-buy-side:benchmark}{tracking error} by trading there; derivatives and fund flows settle on it; the
auction offers the deepest liquidity of the day at no spread; and volume that
moves there thins the intraday market, pushing more volume there. Passive
growth makes each of these stronger every year.
\iqlookfor{benchmarking as the root cause.}
\end{solution}

\begin{solution}{iq:m1:auctions:3}
Sort buy limits descending and sell limits ascending, with market orders at
the front of each; walk the distinct prices once with running cumulative
demand and supply (two pointers), tracking the best executable volume and,
among ties, the smallest surplus: $O(n\log n)$ for the sorts, $O(n)$ for the
sweep. Edge cases: no cross; only market orders on both sides (no limit
price: use the reference); ties over a range of prices; zero-surplus ranges;
quantities that overflow 32 bits; allocation at the clearing price when the
surplus side must be pro-rated or time-prioritised; self-match prevention.
\iqlookfor{the single sweep with running sums, and the market-orders-only
case.}
\end{solution}

\begin{solution}{iq:m1:auctions:4}
The imbalance is more than half an average auction: the \omterm{def:m1:auctions:imbalance}{indicative price}
will be well above the last trade. Offsetting sell interest arrives ---
liquidity providers, and holders with \omterm{def:m1:auctions:moc}{limit-on-close orders} --- shrinking
the imbalance and pulling the \omterm{def:m1:auctions:imbalance}{indicative price} back toward the market, but
not all the way. The continuous price is dragged upward in the last minutes
as those who sell in the auction hedge by buying now, and as others buy ahead
of the close to sell into it. Expect partial reversal at the next open.
\iqlookfor{both the auction dynamics and the spill-over to the continuous
book.}
\end{solution}

\begin{solution}{iq:m1:auctions:5}
Target: \omterm{def:m1:auctions:moc}{closing price} minus the mid at prediction time (in basis points or
spreads), or the next morning's reversal. Features: signed imbalance scaled
by average auction volume and by displayed depth, paired quantity, \omterm{def:m1:auctions:imbalance}{indicative
price} relative to the mid, their changes between messages, time to close,
index events, volatility. Trap: the imbalance feed and the continuous quotes
come with different latencies, and features built from messages published
\emph{after} the decision time leak the answer; evaluation must also account
for the strategy's own order changing the imbalance it conditions on.
\iqlookfor{look-ahead through timestamps, and self-impact.}
\end{solution}

\begin{solution}{iq:m1:auctions:6}
Gains: slower liquidity providers, who can no longer be sniped within a
batch interval, since all orders arriving in the same 100 milliseconds
compete on price, not on time; investors, if the saving is passed on as
narrower spreads. Loses: firms whose edge is pure speed; possibly the venue,
if takers dislike waiting and leave for continuous competitors. Spreads
should narrow where sniping costs were a large part of them and widen nowhere;
but with a single venue adopting it, arbitrage between it and continuous
venues reintroduces a race at each batch boundary, and the benefit depends on
what share of the market moves.
\iqlookfor{price competition replacing time competition, and the
fragmentation caveat.}
\end{solution}
